% MG01. Insert after exc:aw and its elliptic interpretation;
% before exc:codigo. Requires the established L_1 and A_W chart.
\subsubsection{Dual neutrality and orientation of the first joint boundary}
\label{mg01:frontera}

The prolongation boundary joins the three regional histories through
a relation preserving the axis, lateral aggregate, saturation, and
column occupancy. In the established row-vector chart, the final
six-position blocks of the words of length thirty are
\[
 u_4^\pi=110221,\qquad u_4^e=102011,\qquad
 u_4^\varphi=010200.
\]
The superscripts retain the three regions of
\eqref{eq:palabras-regionales}; each word is interpreted as an
element of \(\mathbb F_3^6\), with integer representatives in
\(\{0,1,2\}\). The boundary relation transports the axis by
\begin{equation}
 z:=u_5^\varphi=u_4^e=102011,
 \qquad a:=u_4^\varphi A_WL_1^3=111101.
 \label{mg01:eje-contraste}
\end{equation}
Here \(a\) is the oriented contrast between the two lateral channels
and the axis. The corresponding lateral aggregate is
\begin{equation}
 x+y=210112,\qquad
 q:=x+y+z=012120,\qquad a=x+y-z=111101
 \quad\text{in }\mathbb F_3^6,
 \label{mg01:agregados}
\end{equation}
where \(x=u_5^\pi\) and \(y=u_5^e\). The sum and the contrast have
distinct dual readings: \(qA_W=012000\) and \(aA_W=120210\).

Let \(B\in M_{3,6}(\mathbb F_3)\) be the matrix with rows \(x,y,z\).
Its canonical lift defines, column by column,
\[
 c_j=\left\lfloor\frac{x_j+y_j+z_j}{3}\right\rfloor,
 \qquad
 o_j=\mathbf1_{x_j\ne0}+\mathbf1_{y_j\ne0}+\mathbf1_{z_j\ne0}.
\]
% BEGIN R28_BOUNDARY_OCCUPANCY_INPUT
% BEGIN R28_BOUNDARY_OCCUPANCY
% Insert in MG01 after the definition of c_j,o_j, before their values.
The law for the boundary at depth six fixes \(c_j=1\). With the axis \(z\) and the
aggregate \(q\) already determined, the integer lateral sums are
\[
 (x_j+y_j)_{j=1}^6=(2,4,3,4,4,2).
\]
The minimum-occupancy rule selects, among decompositions with
\(x_j,y_j\in\{0,1,2\}\), the columns with the fewest nonzero entries,
including \(z_j\). For sum two, the extreme pairs \((0,2),(2,0)\)
occupy one lateral position; for sum three, the pairs \((1,2),(2,1)\)
occupy two; for sum four, only \((2,2)\) is possible. Including the
axis gives \(o=(2,2,3,2,3,2)\). At this boundary, with
\(s=[x+y]_3\), the same rule is expressed by
\[
 o_j=2+\mathbf1_{\{z_j\ne0,\ s_j\ne2\}}.
\]
% END R28_BOUNDARY_OCCUPANCY

% END R28_BOUNDARY_OCCUPANCY_INPUT

At this boundary, saturation and occupancy are
\begin{equation}
 c=111111,\qquad o=223232.
 \label{mg01:saturacion}
\end{equation}
Thus the integer equality in each column is
\(x_j+y_j+z_j=q_j+3\). The carry \(c\) records the integer lift
of this boundary; its domain consists of the six columns under consideration.

\begin{proposition}[Eight distributions of a saturated boundary]
\label{mg01:ocho}
Conditions \eqref{mg01:eje-contraste}--\eqref{mg01:saturacion}
leave exactly eight matrices. Their column alternatives are
\begin{equation}
\begin{array}{c|c|c|c}
 j&z_j&x_j+y_j\ \text{in }\mathbb Z&(x_j,y_j)\\\hline
 1&1&2&(0,2),(2,0)\\
 2&0&4&(2,2)\\
 3&2&3&(1,2),(2,1)\\
 4&0&4&(2,2)\\
 5&1&4&(2,2)\\
 6&1&2&(0,2),(2,0)
\end{array}
\label{mg01:tabla-columnas}
\end{equation}
\end{proposition}
\begin{proof}
Axis transport fixes the third row. The integer equality
\(x_j+y_j=q_j+3-z_j\), together with \(o_j\), gives the alternatives
in the table. The first, third, and sixth columns each admit two
independent choices; the other three admit one. Their cardinalities
have product \(2\cdot1\cdot2\cdot1\cdot1\cdot2=8\), and each
choice preserves every stated equality.
\end{proof}

Neutrality is evaluated in the dual chart. With
\(\mathbf1=(1,1,1,1,1,1)^{\mathsf T}\), the matrix in
\eqref{exc:aw} determines
\begin{equation}
 \omega=A_W\mathbf1=(2,1,1,1,1,1)^{\mathsf T},
 \qquad A_W\omega=-\mathbf1.
 \label{mg01:unidad-dual}
\end{equation}
The dual closure condition at the boundary is
\begin{equation}
 (BA_W)\mathbf1=B\omega=0\quad\text{in }\mathbb F_3^3.
 \label{mg01:neutralidad}
\end{equation}
This condition acts on the collection of saturated boundaries. The quadratic
relation for \(A_W\) determines transport of the dual unit; neutrality
specifies which boundaries satisfy the corresponding closure.

\begin{proposition}[Dual pair and oriented boundary]
\label{mg01:dos-orientacion}
Neutrality \eqref{mg01:neutralidad} leaves exactly
\begin{equation}
 B_-=
 \begin{pmatrix}0&2&1&2&2&2\\2&2&2&2&2&0\\1&0&2&0&1&1\end{pmatrix},
 \qquad
 B_+=
 \begin{pmatrix}2&2&2&2&2&0\\0&2&1&2&2&2\\1&0&2&0&1&1\end{pmatrix}.
 \label{mg01:pareja}
\end{equation}
The involution interchanging closure and propagation permutes \(B_-\)
and \(B_+\). In the half-open chart with the APP orientation fixed
for the channels, the canonical boundary is \(R_{36}=B_+\).
\end{proposition}
\begin{proof}
Write the three binary choices in the first row as
\[
 x=(2\varepsilon_1,2,1+\varepsilon_3,2,2,2\varepsilon_6),
 \qquad \varepsilon_1,\varepsilon_3,\varepsilon_6\in\{0,1\}.
\]
The condition \(x\omega=0\) is equivalent to
\[
 1+\varepsilon_1+\varepsilon_3-\varepsilon_6\equiv0\pmod3.
\]
The possibilities are \((0,0,1)\) and \((1,1,0)\). The axis row
satisfies \(z\omega=0\); the lateral aggregate also has zero dual
evaluation. Consequently, \(y\omega=0\) in both cases. This gives
exactly the two matrices in \eqref{mg01:pareja}.

The final choice applies the APP orientation of the half-open chart
to this mirror pair. In that chart, the ordered representations are
\[
 B_-\longmapsto(789,048,894),\qquad
 B_+\longmapsto(793,045,894).
\]
The representative of the fixed orientation is \(B_+\). This selection
preserves the order of the three regions and the boundary convention
of their cylinders. The algebraic reduction from eight matrices to two
and the choice of oriented representative thus constitute two stages
of the same construction.
\end{proof}

The visible charges distinguish these stages precisely:
\begin{equation}
 B_-\mathbf1=(0,1,2)^{\mathsf T},\qquad
 B_+\mathbf1=(1,0,2)^{\mathsf T},\qquad
 B_\pm\omega=0.
 \label{mg01:cargas}
\end{equation}
Visible charge and dual neutrality are evaluations on distinct vectors.
The selected boundary retains three rows of six trits, with ambient
indices \((726,215,301)\), and supplies the junction state \(R_{36}\)
for joint prolongation. Its memory includes the transported axis,
saturation, column incidence, and the orientation distinguishing the
sheets of the pair.
